\documentclass[10pt]{article}
\usepackage[margin=20mm]{geometry}
\usepackage{amsmath,amssymb,mathtools,fontspec,unicode-math,xurl,hyperref}
\IfFontExistsTF{DejaVu Serif}{\setmainfont{DejaVu Serif}}{\setmainfont{Latin Modern Roman}}
\setmathfont{latinmodern-math.otf}
\setlength{\parindent}{0pt}\setlength{\parskip}{6pt}\sloppy
\title{A modular slice obstruction and local corrections
in moduli of Langlands parameters}
\author{AI-assisted technical review}
\date{}
\begin{document}
\maketitle
Abstract: The principal-triple slice used in Appendix B is not an isomorphism under the stated characteristic hypothesis. Explicit examples show failure even beyond the Coxeter number and under the paper’s sufficient banal-prime condition. This invalidates the given slice-based derivation of a determinant identity used in the unobstructedness argument, although the identity itself survives the examples. Four further findings have local repairs: an inertia-cohomology assertion, a finite-image adjustment, an exact sequence on the wrong site, and an inverted prime-set description. The assessment is Major for the proof dependency; the review does not establish that the main moduli-space conclusions are false. Independent scientific acceptance remains pending.\par
\section*{Identity, versions and scope}
The primary source is J.-F. Dat, D. Helm, R. Kurinczuk and G. Moss, Moduli of Langlands parameters, J. Eur. Math. Soc. 27 (2025), 1827–1927, DOI 10.4171/JEMS/1599 [main]. The 101-page publisher PDF records revision on 8 August 2024. The publisher records publication on 11 March 2025. The supplied arXiv:2009.06708v3 was posted on 29 February 2024; it is a distinct, earlier author version, and no equivalence or full collation is asserted.\par
All primary PDF text pages were read. References below use printed journal pages, with one-based PDF pages supplied for each finding; the conversion is printed page \(=\) PDF page \(+1826\). The checks track the tame construction, wild reduction, continuity, connectedness, unobstructed points, affine quotient and appendices. Selected source images, including every page carrying a reported defect, were actually viewed. External theorems and large classification tables were not all independently reproved. This is an initial mathematical review, not independent acceptance.\par
\section*{The principal-triple slice in Appendix B}
Problem 1 (Major).  Appendix B.4, pp. 1923–1924 (PDF 97–98), asserts that \[\mathfrak g^E\longrightarrow\mathfrak g/\!/\widehat G,\qquad
 X\longmapsto(F+X)\bmod\widehat G\tag{1}\] is an isomorphism when \(\widehat G\) is simple adjoint and the coefficient characteristic \(\ell\) does not divide its Weyl-group order. Here \((F,H,E)\) is the specified principal \(\mathfrak{sl}_2\)-triple with \(H=2\rho^\vee\). This assertion is false. Its failure affects the given derivation of Proposition B.5, used explicitly in Proposition 5.18, p. 1886 (PDF 60).\par
\subsection*{A counterexample within the sufficient prime range}
Work over an algebraically closed field \(L\) of characteristic \(7\), with \(\widehat G=\operatorname{PGL}_5\). Since \(5\) is invertible, identify \(\operatorname{Lie}(\operatorname{PGL}_5)\) with \(\mathfrak{sl}_5\). Write \(E_{ij}\) for matrix units and put \[E=E_{12}+E_{23}+E_{34}+E_{45},\quad
 F=4E_{21}+6E_{32}+6E_{43}+4E_{54},\quad
 H=\operatorname{diag}(4,2,0,-2,-4).\] Direct multiplication gives \([E,F]=H\), \([H,E]=2E\), and \([H,F]=-2F\). Thus these are the matrices prescribed by the principal triple. The centralizer of the regular Jordan block \(E\) in \(\mathfrak{sl}_5\) is \[\mathfrak g^E=L E\oplus L E^2\oplus L E^3\oplus L E^4 .\] For \(Y=F+xE+yE^2+zE^3+wE^4\), the only terms contributing to the trace of \(Y^3\) are the three cyclic arrangements of \(F,F,yE^2\). Consequently \[\operatorname{tr}(Y^3)=3y\,\operatorname{tr}(F^2E^2)
         =3y(4\cdot6+6\cdot6+6\cdot4)=252y=0\quad\text{in }L.       \tag{2}\] On the other hand, \(D=\operatorname{diag}(1,1,-2,0,0)\) is traceless and \(\operatorname{tr}(D^3)=-6=1\) in \(L\). The conjugation-invariant function \(\operatorname{tr}(A^3)\) therefore vanishes on the entire proposed slice but is nonzero on \(\mathfrak g/\!/\widehat G\). Map (1) is not surjective.\par
This example satisfies \(\ell\nmid |W|\), since \(|W|=120\), and \(\ell>h=5\). It also persists under the paper’s sufficient prime condition: take residue cardinality \(q=3\). The powers of \(3\) modulo \(7\) are \(3,2,6,4,5,1\), so its order is \(6>h\) and \(\chi_{\widehat G,1}^{*}(3)\ne0\) in \(L\). Merely imposing the stronger condition used for Theorem 5.7 does not restore this particular slice.\par
\subsection*{Two further exact checks}
For \(\operatorname{PGL}_4\) in characteristic \(5\), take the analogous matrices \[E=E_{12}+E_{23}+E_{34},\qquad
 F=3E_{21}+4E_{32}+3E_{43}.\] Here \(5\nmid24\) and \(5>h=4\). Nevertheless \(\operatorname{tr}((F+xE+yE^2+zE^3)^2)=20x=0\), while \(\operatorname{tr}(\operatorname{diag}(1,-1,0,0)^2)=2\). An independent exact determinant calculation gives \[\det(TI-F-xE-yE^2-zE^3)
 =T^4-10xT^2-24yT+9x^2-36z.\] The missing \(T^2\) coefficient modulo \(5\) confirms the obstruction.\par
The preceding claim in B.4 that \(H=2\rho^\vee\) is always regular semisimple under \(\ell\nmid|W|\) also needs correction. In type \(C_2\) and characteristic \(3\), use the symplectic form with antidiagonal entries \((1,-1,1,-1)\). The same \(4\)-by-\(4\) triple has \[H=\operatorname{diag}(0,1,-1,0),\qquad F=E_{32}.\] Since \(|W(C_2)|=8\), this characteristic is permitted. The highest root \(2e_1\) vanishes on \(H\), so \(H\) is not regular. Moreover, \(\mathfrak g^E=L E\oplus L E^3\) and \[\det(TI-F-xE-yE^3)=T^2(T^2-x),\] so the proposed slice has identically zero determinant.\par
\subsection*{Dependence, counterevidence and required repair}
The actual Veldkamp reference [veldkamp], pp. 232, 234–236, uses its specified Chevalley sums \(X_+\) and \(X_-\) and defines the affine space \(X_-+\mathfrak g^{X_+}\). This does not by itself identify that space with the principal-triple expression in (1). The source’s precise choice matters; a name-level invocation of a section theorem does not justify the replacement.\par
The failure of (1) is not a counterexample to Proposition B.5’s final determinant identity. For example, for \(\operatorname{PGL}_n\) in the examples above, the basis \(E,\ldots,E^{n-1}\) has cocharacter weights \(2,4,\ldots,2n-2\). For the identity automorphism, its determinant is directly \[\prod_{j=1}^{n-1}(t^{2j+2}-1)=\chi_{\operatorname{PGL}_n,1}(t^2),\] as B.5 requires. The corresponding type \(C_2\) weights are \(2,6\) and also give the expected formula.\par
A repair must provide a valid characteristic-sensitive proof of the graded, automorphism-equivariant centralizer identity needed for B.5, or a correctly chosen section with all those properties. The formulas just checked prove the identity only in the stated examples. They do not certify a replacement proof for every root system and outer automorphism. A blanket change to \(\ell>h\) is inadequate by (2). The Major assessment concerns this substantive dependency in the unobstructedness argument, not an established failure of flatness, connectedness, or all conclusions of Theorem 5.7.\par
\section*{Four local corrections}
Problem 2 (Minor: the cohomology group must be restricted).  In the proof of Corollary 4.20(1)(b), p. 1869 (PDF 43), the assertion \(H^1(I_F,T_{\phi^\ell}(L))=1\) does not follow from the action factoring through a cyclic \(\ell\)-group.\par
Take residue characteristic \(p=2\), coefficient characteristic \(\ell=3\), \(T=\mathbf G_m\), \(L=\overline{\mathbb F}_3\), and trivial action. A tame quotient \(I_F\twoheadrightarrow C_5\) and an embedding \(C_5\hookrightarrow\mu_5(L)\) give a nontrivial continuous cocycle. All coboundaries for trivial action are trivial, so this \(H^1\) is nonzero.\par
For the fixed-component lifting at issue, the obstruction factors through the finite cyclic \(\ell\)-group \(C\) acting on the full centralizer. The needed group is \(H^1(C,T(L))\), which vanishes because multiplication by \(|C|\) is an automorphism of \(T(L)\). Equivalently one can explicitly restrict to cocycles through \(I_F/I_F^\ell\). This repairs the local use; it does not disprove Corollary 4.20. The quotient notation used on the following page is compatible with this correction.\par
There is a related clarification in Lemma 4.28, p. 1876 (PDF 50): the fixed-point argument may use an abstract infinite cyclic group generating the same finite action, for which \(H^1\) is the group of coinvariants. The literal open subgroup of full inertia is not procyclic. This observation is not used here to claim a counterexample to Lemma 4.28.\par
Problem 3 (Minor: the finite-image adjustment needs an inverse).  In Lemma 4.24(2), p. 1873 (PDF 47), the text chooses \(t^r={}^{L}\phi(\operatorname{Fr})^r\) and claims that \(\phi^t\) has finite image, although the preceding definition multiplies its Frobenius value by \(t\).\par
For \(\widehat G=\mathbf G_m\) with trivial action, \(\phi(s)=1\) and \(\phi(\operatorname{Fr})=2\), each such \(t\) is \(2\zeta\) for a root of unity \(\zeta\), and the new Frobenius value is \(4\zeta\), still of infinite order. Thus the printed step fails even in rank one.\par
For a general repair, write \(g={}^{L}\phi(\operatorname{Fr})\) and let \(U\) be the relevant inertia-fixed torus normalized by \(g\). After increasing \(r\), arrange \(g^r\in(U^{\operatorname{Ad}g})^\circ\). Choose \(t\) in this torus with \(t^r=g^{-r}\). Then \(tg=gt\) and \((tg)^r=1\). The choice remains in the connected inertia centralizer, as required by the earlier deformation argument. Taking an arbitrary root without ensuring commutation would leave a gap. The corrected construction preserves the intended existence conclusion.\par
Problem 4 (Minor: the central-isogeny sequence requires the fppf site).  Lemma A.11(i), p. 1918 (PDF 92), describes its fixed-point/cohomology sequence as exact on the big étale site, although the central-isogeny kernel need not be smooth.\par
Take \(R=\mathbb Z[1/3]\), \(\Gamma=C_3\) acting trivially, and the central isogeny \[1\longrightarrow\mu_2\longrightarrow\operatorname{SL}_2\longrightarrow\operatorname{PGL}_2
 \longrightarrow1.\] These satisfy the printed assumptions. Since \(H^1(C_3,\mu_2)=1\), the displayed étale exactness would force \(\operatorname{SL}_2\to\operatorname{PGL}_2\) to be an étale-sheaf epimorphism. Over \(A=\mathbb F_2[u,u^{-1}]\), however, the class of \(\operatorname{diag}(u,1)\) can have a determinant-one representative only after adjoining a scalar \(a\) satisfying \(a^2u=1\). At the generic point this requires a purely inseparable square root, which no étale extension can supply.\par
The central-isogeny argument should first be made in the fppf topology. One can then return to the étale topology for the finite étale component-group objects. The invertibility of \(|\Gamma|\) is not the invertibility of the isogeny degree. The example refutes the stated site-exactness, not the final finiteness conclusion of A.11 or A.12.\par
Problem 5 (Minor: the introductory prime set is reversed).  On p. 1836 (PDF 10), immediately before and after Theorem 1.9, the description of \(M\) uses all \(^{L}G\)-banal primes rather than the non-banal primes.\par
The preceding page makes the banal set cofinite, so no positive integer can be divisible by every such prime. For the trivial-action torus and \(q=2\), the tame parameter scheme is just \(\mathbf G_m\), and every \(\ell\ne2\) is banal; the literal product is already impossible in this example. Proposition 6.2 and Remark 6.3, p. 1899 (PDF 73), give the intended finite exclusions: invert \(N_{\widehat G}\) and the non-banal primes, with the stated refinement in the unramified case. This is counterevidence against interpreting the introductory wording as a defect in the formal GIT theorem.\par
\section*{Positive checks and proof coverage}
\subsection*{Tame geometry and wild reduction}
For the trivial-action torus, the parameter ring is explicitly \[\mathbb Z[1/p][a^{\pm1},b^{\pm1}]/(b^{q-1}-1).\] It is flat and a complete intersection; its total space is reduced, while its characteristic-\(\ell\) fiber is nonreduced when \(\ell\mid(q-1)\). This directly tests the distinction maintained in Corollary 2.5 and Proposition 2.8. The tame proof allows nonreduced centralizers and uses the reduced strata for its dimension argument.\par
The wild restriction uses finite \(p\)-groups after \(p\) has been inverted. The averaging arguments in Appendix A therefore apply to the actual group orders. Theorem 3.4 guarantees preservation of a Borel pair; the stronger pinning statement has a separate smooth-center hypothesis in Theorem 3.12. The proof chain through torus-valued \(H^2\), finite-image lifts, valuation quotients and finite field extensions was traced with these hypotheses retained. The Bruhat–Tits and reductive fixed-point inputs remain external dependencies.\par
\subsection*{Continuity and connectedness}
The separatedness argument in Corollary 2.11 can be checked by a Noetherian ascending chain: if \(a=\ell^n a_n\ne0\) for every \(n\), stability of \((a_n)\) forces \(\ell\) to be invertible. In the situation at hand, an integral point in characteristic \(\ell\) rules that out. This separatedness makes the closed-subgroup descent of the extended \(\mathrm{GL}_n\) parameter valid: its defining equations vanish modulo every \(\ell^n\) and hence vanish.\par
For Theorem 4.8, a connected special fiber together with separatedness does imply connectedness of the affine total space. An idempotent reducing to zero satisfies \(e=e^n\in\ell^n R\) for every \(n\); hence \(e=0\). The same argument applied to \(1-e\) covers the other reduction. The canonical comparison in Theorem 4.18 uses the jointly faithful local completions to identify the two invariant subrings; it does not require their quotient to be a finitely generated \(\mathbb Z[1/p]\)-module.\par
General Borel-preserving connectedness remains explicitly conjectural. The proved pinning-preserving and smooth-center cases were not silently enlarged. Likewise, canonical independence of the affine quotient is distinct from the proposed equivalence for stack quotients.\par
\subsection*{Unobstructedness and the quotient}
The Euler-characteristic calculation in Lemma 5.1 yields \[\dim T_x Z^1=\dim\widehat G_L+
 \dim H^0(W_F^0,(\operatorname{Ad}\phi_x)^*\otimes\omega).\] Together with the established pure fiber dimension, this gives the smoothness criterion of Corollary 5.3. The later reductions distinguish good primes, very good primes, \(\ell>h\), and the stronger polynomial exclusions; they are not interchangeable. The classical-group weight calculation and the separate triality treatment were read. The exceptional cases use Lawther–Testerman tables, including private matrix information in the \(E_8\) case; these data were not independently reconstructed in full.\par
In Section 6 the restriction to finite inertia, the universal homeomorphism, the localization argument, the torus quotient and comparison with the Haines variety were traced. Theorem 6.8’s assertion about finite Weyl-type group orders was not separately proved in every root system. No failure of that assertion was established, and it is not entered as a finding.\par
\section*{Supporting checks}
The primary PDF is 997918 bytes, SHA256\par
The matrix determinant check uses exact integer polynomial arithmetic, not floating-point sampling. No author code was executed.\par
Three native helpers in the same conversation performed bounded textual checks: `check\_0221\_concentration` read PDF 11–30 and 84–93; `check\_0221\_waiting\_time` read PDF 31–50 with specified introductory and continuity dependencies; and `check\_0222\_mean\_projection` read PDF 51–83 and 94–98. Their exact auxiliary ranges and qualifications are recorded separately. Helpers did not supply an independent acceptance or a verified model-backend identity.\par
The bounded supplement/code/erratum searches found no article-specific code package or correction. This is a search limitation, not evidence of absence and not a defect. The earlier arXiv version was retained without a claim that the publisher corrections or the present findings transfer to it. Complete reproofs of the external classification, pseudocharacter, building, adequate-moduli and deformation theorems, and a complete replacement for the Appendix B dependency, remain outside this initial review.\par
The workflow records `scientific\_decision=pending` and an empty `human\_reviews` array. Controller-observed UI settings are separate from runtime self-description; neither is presented as an independently verified backend or an independent scientific review.\par
Conclusion.\par
Item | State\par
Supported findings | One Major proof-dependency defect and four Minor corrections. The main moduli-space conclusions are not disproved.\par
Repair verification | Counterexamples and stated local calculations checked; a general replacement centralizer proof remains uncertified.\par
Author adoption | Not checked; no communication sent.\par
Remaining work | Replace the slice argument and complete the characteristic-sensitive centralizer and descent arguments.\par
Scientific decision | Pending; human reviews empty; independent acceptance pending.\par
9 J.-F. Dat, D. Helm, R. Kurinczuk and G. Moss, Moduli of Langlands parameters, J. Eur. Math. Soc. 27 (2025), 1827–1927, doi:10.4171/JEMS/1599 (https://doi.org/10.4171/JEMS/1599). Publisher PDF, revision 8 August 2024. The same authors, arXiv:2009.06708v3, 29 February 2024. Retained author version; equivalence to [main] not verified. F. D. Veldkamp, The center of the universal enveloping algebra of a Lie algebra in characteristic \(p\), Ann. Sci. École Norm. Sup. (4) 5 (1972), 217–240, doi:10.24033/asens.1225 (https://doi.org/10.24033/asens.1225). Selected source pages 217–218, 231–236 actually viewed.\par
\end{document}
